\setcounter{section}{1}

\Needspace{5\baselineskip}
\hypertarget{ux79bbux6563ux6d41ux5339ux914d}{%
\section{离散流匹配}\label{ux79bbux6563ux6d41ux5339ux914d}}

\Needspace{5\baselineskip}
\hypertarget{ux4e00ux63d0ux51faux80ccux666f-3}{%
\subsection{一、提出背景}\label{ux4e00ux63d0ux51faux80ccux666f-3}}

目前的生成模型（如Stable Diffusion）在连续空间（像素值是连续的实数）非常成功。但文本是离散的（token ID是整数，如105, 2099），离散空间没有``微小变化''的概念（105变成105.1没有意义），因此直接应用连续扩散模型很困难。

\Needspace{5\baselineskip}
\hypertarget{ux4e8cux6838ux5fc3ux601dux60f3-3}{%
\subsection{二、核心思想}\label{ux4e8cux6838ux5fc3ux601dux60f3-3}}

连续流匹配通过定义一个向量场v\_t(x)，描述样本点x随时间t如何在空间中移动，微分方程为dx/dt=v\_t(x)。在离散空间，我们不能对离散跳变的位置x求导，但概率分布p\_t(x)是随时间连续变化的，我们可以通过连续时间马尔可夫链（CTMC）来建模。

作者发现，如果我们关注从噪声分布p平滑过渡到真实数据分布q的``路径''，定义``概率速度''来描述x的概率分布如何从p转移向q的状态，那么离散版本的公式和连续版本的公式在数学形式上是统一的。这使得很多在连续扩散中好用的技巧（如调度器设计、去噪参数化）可以直接迁移过来。

\Needspace{5\baselineskip}
\hypertarget{ux4e09ux6570ux5b66ux8868ux8fbe-3}{%
\subsection{三、数学表达}\label{ux4e09ux6570ux5b66ux8868ux8fbe-3}}

注意：以下表达式中样本可以取的值的数量（样本空间的元素数）是有限的，概率的值是离散的，p指的是取该样本的实际概率，而不是概率密度函数。

\Needspace{4\baselineskip}
\begin{enumerate}
\def\labelenumi{\arabic{enumi}.}
\tightlist
\item
  概率路径
\end{enumerate}

背景：生成模型的目标是将一个简单的随机噪声分布 \(p_0\)（比如全是掩码 mask 或随机乱码）变成真实的数据分布 \(p_1\)（比如通顺的句子）。

\Needspace{5\baselineskip}
插值 \(p_t\)：我们需要定义一条路径时间 \(t\)（从 \(0\) 到 \(1\)）变化的路径 \(p_t\)：

\begin{itemize}
\tightlist
\item
  当 \(t=0\) 时，\(p_0\) 完全是噪声。
\item
  当 \(t=1\) 时，\(p_1\) 完全是真实数据。
\item
  这条路径是在定义中间过程，在 \((0,1)\) 时，概率分布 \(p_t\) 是如何计算的。
\end{itemize}

\[
p_t(x)=\sum_{x_0,x_1\in\mathcal{D}}p_t(x\mid x_0,x_1)\pi(x_0,x_1)
\]

这个公式描述了整体路径是如何由无数个``个体路径''组成的。
其中，pi(x\_0,x\_1)表示``起始状态为x\_0且终止状态为x\_1的概率''；p\_t(x\textbar x\_0,x\_1)表示``在起始状态为x\_0且终止状态为x\_1的前提下，特定样本x在t时刻出现的概率''（这里x,x\_0,x\_1分别是三个不同样本）。所以这条表达式实际上是全概率公式，通过对不同的(x\_0,x\_1)下在t时刻得到向量x的概率求和，求解出了t时刻采样得到x的总概率。

\Needspace{5\baselineskip}
其中的``个体路径''p\_t(x\textbar x\_0,x\_1)又可以表示为m个概率分布的凸组合：

\[
p_t(x^i\mid x_0,x_1)=\sum_{j=1}^m\kappa_t^j w^j(x^i\mid x_0,x_1)
\]

\Needspace{5\baselineskip}
左边的 \(p_t(x^i\mid x_0,x_1)\) 表示：

\begin{itemize}
\tightlist
\item
  这是目标条件概率。
\item
  它回答了这样一个问题：``如果我们已知起点是 \(x_0\)（比如掩码 \texttt{{[}MASK{]}}），终点是 \(x_1\)（比如单词``猫''），那么在中间某个时刻 \(t\)（比如进度 \(50\%\) 时），第 \(i\) 个位置的 token 取某个值的概率是多少？''
\end{itemize}

\Needspace{5\baselineskip}
右边的 \(\sum_{j=1}^m\) 表示：

\begin{itemize}
\tightlist
\item
  这是一个加权求和。这表明最终的概率分布是由 \(m\) 个基础分布混合而成的。
\end{itemize}

\Needspace{5\baselineskip}
基础分布（Basis Distribution）记为 \(w^j(x^i\mid x_0,x_1)\)，含义如下：

\begin{itemize}
\tightlist
\item
  这是基本的构造块。
\item
\Needspace{5\baselineskip}
  每一个 \(w^j\) 代表一种简单的概率行为。例如：

  \begin{itemize}
  \tightlist
  \item
    \(w^1\) 可能代表``保持为起点 \(x_0^i\)''。
  \item
    \(w^2\) 可能代表``变成终点 \(x_1^i\)''。
  \item
    \(w^3\) 可能代表``变成完全随机的噪声''。
  \end{itemize}
\end{itemize}

因此，该公式定义的是概率的混合：在时刻 \(t\)，这个位置有 \(\kappa_t^1\) 的概率服从分布 \(w^1\)，有 \(\kappa_t^2\) 的概率服从分布 \(w^2\)，以此类推。

\Needspace{5\baselineskip}
对于点x\_i，最终的概率分布是这m个基础分布的加权（凸组合）。若m=2：

\begin{enumerate}
\def\labelenumi{\arabic{enumi}.}
\tightlist
\item
  \(w^1\) 是``保持起点噪声 \(x_0\)''的概率（\(\delta_{x_0}\)）。
\item
  \(w^2\) 是``变成终点数据 \(x_1\)''的概率（\(\delta_{x_1}\)）。
\end{enumerate}

\Needspace{5\baselineskip}
那么公式就变成了：

\[
p_t(x^i\mid x_0,x_1)=(1-\kappa_t)\delta_{x_0^i}+\kappa_t\delta_{x_1^i}
\]

\textbf{物理意义解释}：这句话的意思是，在时间 \(t\)，这个位置的词有 \(\kappa_t\) 的概率已经变成了最终的词 \(x_1\)，有 \(1-\kappa_t\) 的概率还停留在初始的噪声词 \(x_0\)。随着时间 \(t\) 推移，\(\kappa_t\) 逐渐增大，整个序列就慢慢从噪声``流''向了真实数据。

\Needspace{5\baselineskip}
相当于m=2是``在w\_1、w\_2间的线段上取点''，m=3就变成``在w\_1、w\_2、w\_3组成的三角形内取点''，以此类推。几何解释：

假设我们的词表里只有 3 个词：猫、狗、鸟，即 \(d=3\)。

\begin{itemize}
\tightlist
\item
  \textbf{离散状态}：任何一个确定的词（比如``猫''）都可以表示为一个 One-hot 向量：\([1,0,0]\)。
\item
  \textbf{概率分布}：模型输出的不是确定的词，而是一个概率分布，比如 \([0.7,0.2,0.1]\)（\(70\%\) 是猫，\(20\%\) 是狗\ldots\ldots）。
\item
  \textbf{单纯形（Simplex）}：满足``所有元素非负且和为 \(1\)''的向量集合，在几何上构成了一个单纯形。

  \begin{itemize}
  \tightlist
  \item
    对于 3 个词，单纯形是一个三维空间中的等边三角形。
  \item
    三角形的三个顶点分别是 \([1,0,0]\)、\([0,1,0]\)、\([0,0,1]\)，代表确定的词。
  \item
    三角形内部的任何一点都代表一个中间的概率状态。
  \end{itemize}
\end{itemize}

\[
p_t(x^i\mid x_0,x_1)=\sum_{j=1}^{m}\kappa_t^{i,j}w^j(x^i\mid x_0,x_1)
\]

其中的 \(\sum \kappa=1\) 且 \(\kappa\ge 0\)，这正是几何学中``凸组合''的定义。这意味着，无论怎样设计基础路径 \(w^j\)，只要满足这个公式，生成的概率分布 \(p_t\) 永远不会跑出单纯形之外。它保证了模型产生的总是合法的概率分布。

\Needspace{4\baselineskip}
\begin{enumerate}
\def\labelenumi{\arabic{enumi}.}
\setcounter{enumi}{1}
\tightlist
\item
  生成概率速度
\end{enumerate}

\Needspace{5\baselineskip}
\textbf{连续性方程（The Continuity Equation）} 为了定义离散流匹配，我们使用连续时间马尔可夫链（CTMC）范式。样本 \(X_t\) 在离散状态间跳跃。我们需要预测每个 token 概率变化的\textbf{速率}。连续性方程描述了状态概率的变化率 \(\dot{p}_t(x)\) 与通量（flux）之间的关系：

\[
\dot{p}_t(x)+\mathrm{div}_x(p_tu_t)=0
\]

\Needspace{5\baselineskip}
其中 \(\mathrm{div}_x\) 是散度算子。在离散情况下，散度定义为流出通量减去流入通量：

\[
\mathrm{div}_x(v)=\sum_{z\in\mathcal{D}}[v(z,x)-v(x,z)]
\]

\Needspace{5\baselineskip}
\textbf{定理 2（边际速度）} 给定生成条件概率路径 \(p_t(x\mid x_0,x_1)\) 的条件概率速度 \(u_t^i(x^i,z\mid x_0,x_1)\)，其边际速度 \(u_t\) 由下式给出：

\[
u_t^i(x^i,z)=\sum_{x_0,x_1\in\mathcal{D}}u_t^i(x^i,z\mid x_0,x_1)p_t(x_0,x_1\mid z)
\]

这与连续流匹配中的逻辑一致：边际向量场是条件向量场的后验期望。

\Needspace{5\baselineskip}
\textbf{定理 3（闭式解）} 对于由 Eq. 9 定义的路径（即线性插值路径），边际生成概率速度具有非常简洁的闭式解，且与连续流匹配的形式惊人一致：

\[
u_t^i(x^i,z)=\frac{\dot{\kappa}_t}{1-\kappa_t}[p_{1\mid t}(x^i\mid z)-\delta_z(x^i)]
\]

其中 \(p_{1\mid t}(x^i\mid z)=\mathbb{E}[X_1\mid X_t=z]\) 是\textbf{概率去噪器（Probability Denoiser）}。这意味着速度场可以通过学习一个预测目标 \(X_1\) 分布的神经网络来参数化。

\Needspace{4\baselineskip}
\begin{enumerate}
\def\labelenumi{\arabic{enumi}.}
\setcounter{enumi}{2}
\tightlist
\item
  训练目标
\end{enumerate}

\textbf{定理 3 的威力}：它把复杂的``速度计算''转化为了一个简单的监督学习任务。公式 \(u_t\propto p_{1\mid t}-\delta_z\) 告诉我们：只要训练一个神经网络，输入当前带噪声的句子 \(z\) 和时间 \(t\)，预测原始干净句子 \(X_1\) 的概率分布 \(p_{1\mid t}\)，这个网络就自动定义了正确的``流''速度。

\Needspace{5\baselineskip}
\textbf{训练损失函数}：基于上述推导，训练目标简化为最小化交叉熵损失（Cross Entropy）：

\[
\mathcal{L}(\theta)=-\sum_i\mathbb{E}_{t,data}\log p_{1\mid t}(X_1^i\mid X_t;\theta)
\]

这就是标准的\textbf{去噪（Denoising）}目标：给模型看加了噪声/掩码的数据 \(X_t\)，让它预测原数据 \(X_1\)。这解释了为什么 BERT 风格的 Masked Language Modeling（MLM）实际上是在训练一个流模型。

当然，这与BERT的掩码语言模型（MLM）训练有一个关键区别：掩码率是动态的。在t=0时几乎全被掩盖，在t=1时几乎全是真实数据。模型的输入不只有加噪数据X\_t，而是为(X\_t,t)，意味着它需要通过学习来具备在不同噪声水平（与时间t对应）下恢复原始数据X\_1的能力。

\FloatBarrier
\Needspace{5\baselineskip}
\hypertarget{ux53c2ux8003ux6587ux732e-136}{%
\subsection{参考文献}\label{ux53c2ux8003ux6587ux732e-136}}

\vspace{0.25em}
\begin{itemize}
\tightlist
\item
  Gat, I., Remez, T., Shaul, N., Kreuk, F., Chen, R. T. Q., Synnaeve, G., Adi, Y., \& Lipman, Y. (2024). \href{https://arxiv.org/abs/2407.15595}{Discrete Flow Matching}. arXiv:2407.15595.
\item
  Campbell, A., Yim, J., Barzilay, R., Rainforth, T., \& Jaakkola, T. (2024). \href{https://arxiv.org/abs/2402.04997}{Generative Flows on Discrete State-Spaces: Enabling Multimodal Flows with Applications to Protein Co-Design}. arXiv:2402.04997.
\end{itemize}

\clearpage
